\scanpage{11}
\uline{Remarques} Des \deleted{formules de}\inserted{on pose}
\[
 \underset{\text{(discriminant)}}{\delta(\varphi)}=\Delta\omega^{\unsure{\otimes}{-2}},\qquad\Delta\in\Gamma(S,\underline{\mathrm O}_S) \tag*{(41)}
\]
et les formules de Kramer donnent alors
\[
 \left\{\begin{aligned}
 f'f&=\Delta\,\mathrm{id}_E\\
 ff'&=\Delta\,\mathrm{id}_{E'}
 \end{aligned}\right. \tag*{(42)}
\]
On se gardera de croire que les relations entre $f,f'$ sont symétriques. P.ex. si $f\unsure{^{a}|}$ est de rang $\leq1$ (i.e. les mineurs de rang 2 de $f$ sont nuls) $f'$ est \uline{nul}, et réciproquement, tandis que si $f$ est nul, $f'$ est nul lui-$=$ (et non seulement son rang est $\leq1$)!

Notons que si $\omega$ est multiplié par $\lambda\in\Gamma(S,\underline{\mathrm O}_S^{\unsure{*|\times}})$, ($\varphi$ donc $f$ restant invariant) $\Delta$ et $f'$ sont multipliés l'un et l'autre par $\lambda^2$ -- tandis que le crochet $[\, ,\,]_\varphi$ \inserted{sur $E'$} est multiplié par $\lambda$, et \inserted{le crochet} $[\, ,\,]_{\varphi'}$ sur $E'$, pour \deleted{\unsure{et}} une $\varphi'$ fixée (indépendamment d'une $\varphi$ sur $E$) est multiplié (comme $\omega'=\omega^{\otimes(-1)}$) par $\lambda^{-1}$.

Mais si $\varphi'$ est la forme associée à $f$ via $f'$, comme \inserted{$f'$ donc} $\varphi'$ est \deleted{multiplié}\inserted{remplacé} par $\lambda^2\varphi'$, $[\, ,\,]_{\varphi'}^{\unsure{\omega|}}$ est remplacé par $\lambda[\, ,\,]_{\varphi'}$ : Ainsi $[\, ,\,]_\varphi$ \inserted{sur $E$} et $[\, ,\,]_{\varphi'}$ sur $E'$, associés à un couple

% archive sheet number: 10
